% ==============================================================================
% SECTION 5: EXPERIMENTAL EVALUATION  (Primary author: Anaswara)
%
% Tables/figures are generated by reproduction/run_all.py into paper/tables and
% paper/figures. Every numeric claim in the text marked \resultpending{...}
% must be filled from those files after the cluster runs -- do not invent numbers.
% ==============================================================================
\providecommand{\resultpending}[1]{\textcolor{red}{[#1]}}
\providecommand{\inputifexists}[2]{\IfFileExists{#1}{\input{#1}}{\noindent\resultpending{#2}\par}}
\providecommand{\figifexists}[3]{\IfFileExists{#1}{%
  \begin{figure}[t]\centering\includegraphics[width=#2\linewidth]{#1}\caption{#3}\end{figure}}%
  {\noindent\resultpending{figure \detokenize{#1} not generated yet}\par}}

We evaluate whether MOCAP (i) selects faster plans than budget-unaware and
greedy baselines on the same candidate set, (ii) keeps the realised budget
violation rate near the target $\delta$, (iii) remains reliable when the
assumptions behind its predictions are violated, and (iv) tracks real cloud
spend.

\subsection{Experimental Setup}
\label{sec:exp_setup}

\paragraph{Cluster and software.}
All measurements use Apache Spark~4.2.0 (Python~3.12, Java~17) on
\resultpending{provider, e.g.\ Google Dataproc}: one
\resultpending{master type} master and \resultpending{$k$ $\times$ worker type}
workers (\resultpending{total} executor cores, \resultpending{total} GB executor
memory). The exact configuration of every run is versioned in
\texttt{configs/spark\_env.yaml} and recorded next to each result file.

\paragraph{Data and workload.}
We use data generated by the TPC-DS toolkit at scale factors 10 and 100
(\texttt{scripts/generate\_tpcds.py}). Results are \emph{derived from} TPC-DS
and are not audited TPC-DS results. The main workload is a star-join family
over \texttt{store\_sales} with 3--10 relations, each at three date-range
selectivities, plus a four-dimension group-by query (27 queries in four
families). \resultpending{If run: the official 99 TPC-DS queries that the
candidate generator supports ($n$ of 99); unsupported queries are listed in the
artifact.}

\paragraph{Measurement protocol.}
For each query, the candidate generator of Section~\ref{sec:system_design}
(beam width 40, analysis cap 48) produces the candidate pool, and every
candidate receives a pre-execution prediction. We then \emph{execute every
candidate} \resultpending{10} times after one discarded warm-up run, with
Spark's join reordering disabled and automatic broadcast disabled so that the
executed plan is the hinted plan; each run's executed physical plan is
fingerprinted and checked against the intended plan. Spark's own planner is
additionally run on the original SQL with CBO on/off and AQE on/off, using
Spark's default broadcast threshold. Latency is wall-clock time of a full
execution into a no-op sink; bytes scanned, bytes shuffled and executor time
are read from Spark's monitoring API; monetary cost applies the pricing model
of Section~\ref{sec:system_design} to these measurements. All runs are logged in
one telemetry schema (query, plan, strategy, latency, bytes, executor seconds,
run number, cache state, AQE/CBO flags, cores, memory, seed, timestamp).

Because all strategies choose from the \emph{same} measured pool, a strategy's
outcome for a query is the measured distribution of the plan it selects; the
comparison is therefore paired per query and isolates the selection rule from
candidate generation.

\paragraph{Budgets and calibration.}
Following our audit guidelines, budgets are relative to the Oracle: for a query
whose cheapest measured candidate costs $C^\star$, the budget is
$B = m \cdot C^\star$ with $m \in \{1.0, 1.2, 1.5, 2.0\}$ (default $m=1.5$).
The conformal quantile used for query $q$ is calibrated only on candidates of
the other queries (leave-one-query-out), so no execution of $q$ influences its
own bound. Unless stated otherwise $\delta = 0.05$.
\resultpending{State which conformal score was used: signed $\log(C/\hat C)$
(one-sided) or absolute $|\log(C/\hat C)|$; see \texttt{e1\_run\_info.json}.}

\paragraph{Strategies.}
\textbf{MOCAP} applies the conformal filter $\hat C^{+} \le B$, the Pareto
frontier and the fastest-under-budget rule (Section~\ref{sec:theory}).
\textbf{MOCAP (point)} is the same pipeline filtering on point predictions.
\textbf{Greedy Cheapest} takes the lowest predicted cost.
\textbf{Greedy Fastest} takes the lowest predicted latency among plans with
predicted cost $\le B$. \textbf{Random} picks uniformly (seeded).
\textbf{Spark} denotes Spark's own planner on the original SQL in four CBO/AQE
configurations; \emph{Spark default} is CBO off, AQE on (Spark~4 defaults).
\textbf{Oracle} picks the fastest plan whose \emph{measured} cost is within
budget; it is an upper bound, not a deployable baseline.

\paragraph{Metrics and statistics.}
We report the median latency of the selected plans across queries with a 95\%
percentile-bootstrap confidence interval (2{,}000 resamples), P90/P99 latency
over all runs of the selected plans, the budget-violation rate (fraction of
queries whose selected plan's median measured cost exceeds $B$), the
feasibility rate (fraction of queries for which a strategy found a plan it
believed feasible), Oracle regret $L/L_{\text{Oracle}}-1$, and geometric-mean
speedup over Spark default. Differences are tested with two-sided Wilcoxon
signed-rank tests paired by query, with Holm correction across baselines.

\subsection{E1: Plan Selection Quality}
\label{sec:e1}

\inputifexists{tables/e1_main.tex}{Table e1\_main: run reproduction/run\_all.py}

\resultpending{Summarise the E1 results once the table is generated: MOCAP's median latency vs.\
Greedy and Spark default with CIs; Wilcoxon wins/losses/ties and Holm-adjusted
$p$; budget-violation rate of MOCAP vs.\ MOCAP (point) and Greedy Fastest;
regret to Oracle. Report where MOCAP is slower as well.}

\figifexists{figures/e1_latency_cdf.pdf}{0.6}{Distribution over queries of each
strategy's latency relative to the Oracle (E1).}

\subsection{E2: Budget Sweep}
\label{sec:e2}

\inputifexists{tables/e2_budget.tex}{Table e2\_budget: run reproduction/run\_all.py}

\figifexists{figures/e2_budget_sweep.pdf}{0.9}{Budget-violation rate and median
latency as the budget multiplier $m$ grows (E2).}

\resultpending{Describe the trade-off: violation rate of MOCAP relative to
$\delta$ at each $m$; how point-filter and greedy strategies behave at tight
budgets ($m=1.0, 1.2$); latency cost of the conservative bound; effect of
$\delta \in \{0.01, 0.05, 0.10, 0.20\}$ (full table in \texttt{e2\_summary.csv}).}

\subsection{E3: Planning Overhead}
\label{sec:e3}

Planning overhead and pruning quality are reported in
Section~\ref{sec:system_design}. \resultpending{Relate the median planning time
per query to the median latency saved over Greedy / Spark default in E1.}

\subsection{E4: Ablations}
\label{sec:e4}

\resultpending{Summarise the component ablations (conformal bound vs.\ point
filter, Pareto frontier, decision rule) from \texttt{experiments/ablations.py};
the E1 row ``MOCAP (point)'' is the conformal-bound ablation on real executions.}

\subsection{E5: Where MOCAP Beats Greedy}
\label{sec:e5}

\inputifexists{tables/e5_case_study.tex}{Table e5\_case\_study: run reproduction/run\_all.py}

\resultpending{Name the query family (or families) where MOCAP is significantly
faster than, or violates the budget less often than, Greedy, with a concrete
query from \texttt{e5\_examples.csv} explaining why Greedy's choice failed
(predicted vs.\ measured cost). If no family reaches significance, say so.}

\subsection{Robustness}
\label{sec:robustness}

Each stressed scenario is calibrated on the clean world and tested on the
stressed one: missing statistics, stale statistics (data grown
\resultpending{$4\times$} after \texttt{ANALYZE}), injected estimation errors
(plan-specific factors in $[0.1,1]$ and $[1,10]$), Zipf-skewed foreign keys,
a scale shift (calibrate at SF10, test at SF100), and 2 or 4 concurrent
background jobs. For injected errors we also report the case where calibration
uses the same biased estimator.

\inputifexists{tables/robustness.tex}{Table robustness: run reproduction/run\_all.py}

\figifexists{figures/robustness_violation.pdf}{0.85}{Budget-violation rate per
stress scenario; dashed line marks $\delta$.}

\resultpending{Discuss which shifts preserve coverage and which break it. In our
laptop smoke test, coverage fell when only the test-time estimator was biased
low and under a scale shift, and recovered when calibration saw the same
estimator, i.e.\ conformal guarantees hold under exchangeability, not under
distribution shift. Confirm or refute on the cluster.}

\subsection{Real Cloud Cost}
\label{sec:cloud_cost}

Each experiment batch ran on its own labelled cluster; billing was exported per
label and compared with the provider cost implied by the batch's telemetry.

\figifexists{figures/real_cost.pdf}{0.5}{Telemetry-derived vs.\ billed cost per
batch.}

\resultpending{Spearman correlation between telemetry-derived and billed cost
per batch, and between pre-execution predicted cost and provider cost per plan
(\texttt{cloud\_cost\_report.json}). Note that bills include idle time and
master nodes.}

\subsection{Limitations and Threats to Validity}
\label{sec:limitations}

\begin{itemize}
  \item \textbf{Marginal, not conditional, coverage.} Split conformal
  prediction certifies $\Pr(C \le \hat C^{+}) \ge 1-\delta$ for a plan drawn
  exchangeably with the calibration plans. The plan MOCAP \emph{selects} is not
  such a draw (selection favours plans whose prediction is low), so the
  realised violation rate of selected plans can exceed $\delta$; we therefore
  report it directly rather than assume it.
  \item \textbf{Exchangeability.} Coverage relies on calibration and test
  queries being exchangeable; the robustness study shows that shifts in data
  scale or estimator bias can break it.
  \item \textbf{Candidate coverage.} The generator handles inner equi-join
  queries; queries it cannot rewrite are excluded and listed.
  \item \textbf{Cost model.} Monetary cost is a model of executor time and
  bytes; real bills are hourly per cluster and are compared only by rank
  correlation per batch.
  \item \textbf{Measurement noise and caching.} Medians over repeated runs
  reduce but do not remove noise; warm-cache runs understate cold-start I/O.
  \item \textbf{Data.} Results use TPC-DS-derived data on one cluster shape and
  may not generalise to other workloads or hardware.
\end{itemize}
